% ===========================================================================
% APPENDIX B: IMPLEMENTATION DETAILS
% ===========================================================================
\appendix
\setcounter{section}{1}\renewcommand{\thesection}{\Alph{section}}
\setcounter{subsection}{0}
\renewcommand{\thesubsection}{\thesection.\arabic{subsection}}
\setcounter{table}{0}
\renewcommand{\thetable}{B\arabic{table}}

\section{Implementation Details}\label{app:implementation}

\subsection*{Package Structure}\label{sec:package}

\begin{center}
\small
\begin{tabular}{ll}
\toprule
\texttt{adaptsimex/}\\
\quad\texttt{R/main.R} & main \texttt{adaptive\_simex()} driver\\
\quad\texttt{R/cv\_select.R} & cross-validation for
$(\lambda_{\max},\text{degree})$ (prefix property, 1-SE rule)\\
\quad\texttt{R/variance.R} & analytic sandwich (equation~3.4 of the main text)
and pairs-bootstrap\\
\quad\texttt{R/methods.R} & S3 methods: print, summary, confint, plot\\
\quad\texttt{src/simex\_batch.cpp} & optional Rcpp batching of the
replicate-level IRLS\\
\quad\texttt{vignettes/} & quick start; theory; NHANES case study\\
\quad\texttt{tests/testthat/} & unit tests (attenuation direction,
sandwich vs.\ bootstrap)\\
\bottomrule
\end{tabular}
\end{center}

\subsection*{Main Function and Usage}\label{sec:usage}

\begin{center}
\small
\begin{minipage}{0.93\textwidth}
\texttt{adaptive\_simex(formula, data, me\_vars, sigma\_error,\\
\quad lambda\_max\_seq = NULL, poly\_degree\_seq = 1:3,\\
\quad n\_simex = 100, n\_cv\_folds = 5,\\
\quad se\_method = c("sandwich", "bootstrap"), n\_boot = NULL,\\
\quad parallel = TRUE, n\_cores = NULL, verbose = TRUE)}\\
Returns an object of class \texttt{c("adaptive\_simex", "glm")} with
slots \texttt{coefficients}, \texttt{vcov}, \texttt{se},
\texttt{selected\_lambda\_max}, \texttt{selected\_poly\_degree},
\texttt{cv\_results}, and \texttt{call}.
\end{minipage}
\end{center}

\noindent Here \texttt{me\_vars} names the error-contaminated columns
and \texttt{sigma\_error} their standard deviations;
\texttt{se\_method = "sandwich"} uses \eqref{eq:sandwich}, and
\texttt{"bootstrap"} the pairs-bootstrap recommended at small $n$
(main text, Section~5.7). A typical call:

\begin{center}
\small
\begin{minipage}{0.93\textwidth}
\texttt{fit <- adaptive\_simex(Y \textasciitilde W1 + W2 + X3, data = df,\\
\quad me\_vars = c("W1","W2"), sigma\_error = c(0.3, 0.3),\\
\quad n\_simex = 100, se\_method = "sandwich", parallel = TRUE)}\\
\texttt{summary(fit); confint(fit); plot(fit)}
\end{minipage}
\end{center}

\subsection*{Computational Complexity and Parallelization}\label{sec:complexity}


Three implementation ideas reduce the cost far below naive estimates.
First, all $B$ SIMEX replicates at a grid point are fitted by one
\emph{batched} IRLS sweep: the score and Hessian accumulations are
tensor contractions over a $(B,n,p)$ array, so the $B$ fits share the
linear-algebra machinery. Second, the CV procedure exploits the
\emph{prefix property} of the SIMEX grid: one union-grid curve per
fold (with noise shared across $\lambda$ within a replicate, per
Algorithm~1) serves all candidate $\lambda_{\max}$, reducing the
number of fitted pseudo-samples per fold from
$|\mathcal K_\lambda|\cdot|\Lambda|\,B_{\mathrm{cv}}$ to
$|\Lambda|\,B_{\mathrm{cv}}$. Third, fits along the curve are warm
started from the previous grid point. Measured single-core times for
the full A-SIMEX fit (CV plus final curve plus sandwich) are $0.10$,
$0.28$, $0.95$, and $3.3$\,seconds at $(n{=}200,p)\in
\{5,15,30,50\}$ with $B=100$, $B_{\mathrm{cv}}=25$; standard SIMEX is
about twice as fast, and the corrected-score estimator is one to two
orders of magnitude cheaper but numerically fragile
(main text, Section~5.3). The asymptotic complexity per replication
remains $O(B\,|\Lambda|\,n p^2)$ for the final curve plus
$O(5\,|\Lambda|\,B_{\mathrm{cv}}\,n p^2)$ for the CV, but the
constants are small enough that parallelization over Monte Carlo
repetitions, CV folds, or grid points, all of them embarrassingly
parallel and all supported by the package through the
\texttt{future}/\texttt{furrr} backends, is a matter of convenience
rather than necessity on modern hardware.

% ---------------------------------------------------------------------------
% SECTION 8: DISCUSSION (revised)

